\section{Related Work}

\subsection{Psychophysics, systems neuroscience, and behavioral reverse engineering}

Our methodology is inspired by the broader logic of psychophysics and systems neuroscience, where the goal is not merely to score an observer on isolated stimuli, but to infer the structure of its latent representation through controlled probes and perturbations. Representational similarity analysis emphasizes that structure can be inferred behaviorally even when the underlying mechanism is not directly accessible \cite{kriegeskorte2008rsa}. Goal-driven modeling in visual neuroscience similarly argues that what matters is not one task score in isolation, but whether a model reproduces systematic representational structure and robust invariances \cite{yamins2016dicarlo}. We adopt this stance for VLMs: the paper does not claim to decode the model's internal neural state directly, but instead studies the externalized scene belief implied by its answers.

\subsection{Controlled synthetic worlds for visual reasoning}

CLEVR and CLEVRER established the scientific value of controlled synthetic benchmarks when the aim is to study reasoning rather than ecological realism alone \cite{johnson2017clevr,yi2020clevrer}. By controlling object attributes, compositional structure, and event dynamics, these benchmarks make it possible to localize failure modes much more precisely than open-world datasets typically allow. Ordinary-Bench inherits this experimental logic, but focuses specifically on whether a model's answers imply a reconstructable spatial world. For this purpose, exact geometry, controllable complexity, and repeatable visual perturbations are essential.

\subsection{Spatial benchmarks for VLMs}

Recent work has shown that spatial understanding remains a persistent weakness for VLMs. SpatialVLM targets spatial reasoning more directly \cite{chen2024spatialvlm}, while VSR and What'sUp document failures on visual spatial relations and directional understanding \cite{liu2023vsr,kamath2023whatsup}. These benchmarks are important, but they still largely frame the problem as one of question-level correctness. They ask whether the model can solve a spatial item, not whether all of its answers together define a coherent scene. Ordinary-Bench differs by taking scene reconstructability rather than local correctness as the central criterion.

\subsection{Language priors, hallucination, and evaluation mismatch}

Another line of work argues that benchmark success can be inflated by language priors or evaluation artifacts. BLINK highlights the gap between seeing and perceiving in multimodal models \cite{fu2024blink}; MMStar questions whether current VLM evaluation pipelines measure the intended abilities at all \cite{chen2024mmstar}; POPE studies object hallucination \cite{li2023pope}; and VLInd-Bench explicitly measures the contamination of visual answers by language priors \cite{lee2025vlind}. Ordinary-Bench is aligned with this critique, but pushes it one step further. The issue is not only whether a model can answer correctly without vision, but whether its answers support any scene that could actually be reconstructed.

\subsection{Selective answering and reliability}

Reliable multimodal evaluation should also account for uncertainty and abstention. Reliable VQA argues that abstaining is often better than confidently answering incorrectly \cite{whitehead2022reliablevqa}; Selectively Answering Visual Questions and later consistency-based selective VQA frameworks extend this argument to black-box VLMs \cite{eisenschlos2024selectivevqa,khan2024consistency}. This perspective is highly relevant here because once answers are interpreted as geometric constraints, a confidently wrong answer is not just a single mistake: it can corrupt the entire reconstructed scene. For this reason, Ordinary-Bench is naturally compatible with selective-answering and consistency-oriented protocols.

\subsection{Ordinal embedding and geometric realizability}

The immediate theoretical basis of our evaluation lies in ordinal embedding and geometric realizability. Generalized non-metric multidimensional scaling and later robust ordinal embedding work show that relative comparisons can recover geometric structure even without absolute metric values \cite{agarwal2007gnmds,ma2018robustordinal}. Euclidean distance matrix theory reminds us, however, that a set of ordinal constraints is not automatically geometrically realizable \cite{dokmanic2015edm}. For directional relations, bearing and angle rigidity provide the right language for understanding uniqueness, mirror ambiguity, and the realizability of directional constraints \cite{zhao2019bearingrigidity,chen2021anglerigidity}. Ordinary-Bench does not claim a new general-purpose reconstruction algorithm; instead, it turns these ideas into an evaluation principle for VLMs: if enough accurate ordinal relations can reconstruct a scene, then dense relational probing can reveal whether a model possesses a coherent scene belief at all.

\subsection{How our work differs}

Existing work typically concentrates on one of three questions: whether VLMs can answer spatial questions, whether visual evaluation is contaminated by language priors, or whether ordinal constraints can reconstruct geometry. Ordinary-Bench combines all three. It adopts a psychophysics-inspired black-box observer view, uses controlled synthetic scenes, and treats geometric reconstructability as the end-point of evaluation. In this sense, the benchmark is not simply another spatial QA dataset; it proposes a stricter criterion for vision-language understanding: a model should count as spatially competent only if its many local answers jointly support a coherent, reconstructable, and alignable scene.
